\documentclass[12pt,a4paper]{article}
\usepackage[utf8]{inputenc}
\usepackage{graphicx}
\usepackage{hyperref}
\usepackage{geometry}
\geometry{a4paper, margin=1in}

\title{SIT223-SIT753 Sprint 1 Group Task \\ \large Version Control Report}
\author{Team Leader: Prabhath}
\date{\today}

\begin{document}

\maketitle

\section{Introduction}
This report outlines the version control processes and collaborative efforts undertaken during Sprint 1 of the SIT223-SIT753 Unit Page development. The team utilized Git and GitHub to branch, develop features independently, and merge them into the centralized project repository.

\section{Methodology and Roles}
The project was divided into three primary roles. Each member branched off \texttt{main}, implemented their respective features, and merged their contributions:
\begin{itemize}
    \item \textbf{UI Designer:} Created the \texttt{add-header} branch. Added the Deakin logo, header layout, and corresponding styling.
    \item \textbf{Front-end Developer:} Created the \texttt{add-table} branch. Built the unit details table with specific trimester dates, enrolment modes, and zebra-striping CSS.
    \item \textbf{Team Leader (Prabhath):} Created the \texttt{add-unit-content} branch. Added unit descriptions, fee links, and hurdle requirements. Managed the overarching GitHub repository workflow.
\end{itemize}

\section{Execution and Merging}
All team members successfully checked out from \texttt{main}, applied their respective changes, and merged them back. Any consecutive merge conflicts (such as overlapping CSS styles in \texttt{main.css}) were manually resolved to ensure a stable, cleanly formatted final product. Following the merges, each team member made a final autonomous commit to refine the aesthetics and content.

\section{Repository Insights}

\subsection{Contributors Graph}
The following image displays the Contributions Graph, verifying the commits made by each respective team role over the sprint.

\begin{figure}[htb]
    \centering
    \includegraphics[width=0.8\textwidth]{1.png}
    \caption{GitHub Contributors Graph}
    \label{fig:contributors}
\end{figure}

\subsection{Network Graph}
The Network Graph below visualizes the git history, clearly illustrating where the three parallel branches (\texttt{add-header}, \texttt{add-table}, and \texttt{add-unit-content}) diverged and when they were subsequently merged back into \texttt{main}.

\begin{figure}[htb]
    \centering
    \includegraphics[width=0.8\textwidth]{2.png}
    \caption{GitHub Network Graph}
    \label{fig:network}
\end{figure}

\section{Project Management (Trello)}
The team utilized a Trello board to assign tasks, assign roles, and track the progress of the Sprint 1 requirements.

\begin{figure}[htb]
    \centering
    \includegraphics[width=0.8\textwidth]{3.png}
    \caption{Sprint 1 Trello Board Status}
    \label{fig:trello}
\end{figure}

\end{document}
